\section{Supporting Verified Lemmas}
\label{sec:supporting-lemmas}

The theorem above is the article's main result. We also record a few
closely related lemmas that reuse the same prime and divisibility
foundations. They are included here as supporting results, not as
additional headline claims.

\begin{itemize}
  \item Finite-prefix consequence: a complete finite prime prefix has a
    larger prime (Section~4.1).
  \item Divisor-search bounds: composite inputs have a proper smallest
    divisor below their square root (Section~4.2).
  \item Finite-prefix primality: coprimality plus factor coverage excludes
    all proper divisors (Section~4.3).
  \item B\'ezout product lemmas: a prime divisor of a product is forced
    onto a factor (Section~4.4).
\end{itemize}

\subsection{Corollary: Greater Than a Complete Finite Prefix}
\label{sec:complete-finite-prefix}

A direct corollary of Euclid's construction is that a complete finite
prefix of the primes is never closed. Let $P=[p_1,\dots,p_k]$ be a sorted
finite list that contains every prime up to its largest element $h=p_k$.
Let $q$ be the prime produced by the Euclid construction from $P$. Since
Section~3 proves $q\notin P$, $q$ cannot be at or below $h$: every prime
at or below $h$ is already contained in the complete prefix. Therefore
$q > h$.

\begin{equation*}
\begin{aligned}
P &= [p_1,\dots,p_k],\quad h=p_k
&&\text{[Finite Prime Prefix]} \\
\forall r,\ \operatorname{isPrime}(r)\land r\le h
&\Rightarrow r\in P
&&\text{[Prefix Complete Through }h\text{]} \\
\operatorname{isPrime}(q)
&\land q\notin P
&&\text{[Euclid Construction]} \\
q\le h
&\Rightarrow q\in P
&&\text{[Prefix Completeness]} \\
q\le h
&\Rightarrow q\in P\land q\notin P
&&\text{[Contradiction]}
\end{aligned}
\end{equation*}

\begin{equation*}
\therefore\ q>h
\quad\blacksquare\ \text{[Q.E.D.]}
\end{equation*}

This is the ordered-prefix form of the Euclid theorem: from any complete
finite prefix of the primes, the construction produces a prime beyond that
prefix.

{\raggedright
This corollary is verified in
\href{https://github.com/thiagomata/prime-numbers/blob/master/src/main/scala/v1/chapter5/prime/properties/PrimeProperties.scala}{\texttt{PrimeProperties::newPrimeNotInList}},
\href{https://github.com/thiagomata/prime-numbers/blob/master/src/main/scala/v1/chapter5/prime/properties/PrimeProperties.scala}{\texttt{PrimeProperties::notContainsFromValueNotMatchesAny}},
and
\href{https://github.com/thiagomata/prime-numbers/blob/master/src/main/scala/v1/chapter5/prime/properties/PrimeProperties.scala}{\texttt{PrimeProperties::euclidPrimeGreaterThanHead}}.
\par}

\subsection{Smallest-Divisor Bounds for Composite Numbers}
\label{sec:smallest-divisor-bounds}

The primality test used in Euclid's proof relies on
\texttt{findSmallestDivisor(n, 2)}, which scans candidates from 2 upward
until it finds the smallest divisor of $n$. Two lemmas characterize why
this scan is both correct and efficient.

\textbf{Composite has a divisor below n.} If $n$ is composite and $d$ is
its smallest non-trivial divisor, then $d < n$. This is immediate from the
definition of composite --- there exists a proper divisor --- but must be
proved against the \texttt{findSmallestDivisor} algorithm, which scans
upward until a divisor is found or $n$ itself is reached.

\begin{equation*}
\begin{aligned}
n > 1 \;\land\; \neg \operatorname{isPrime}(n) &\Rightarrow \\
\exists\, d = \operatorname{findSmallestDivisor}(n, 2) &: 2 \leq d < n \;\land\; \operatorname{Calc.mod}(n, d) = 0
\end{aligned}
\end{equation*}

\textbf{Proof.} Since $n$ is composite, it has a proper divisor
$e\in[2,n)$. The minimality specification from Section~2.2 gives
$2\leq d\leq e<n$ and $\operatorname{mod}(n,d)=0$.

\begin{equation*}
\therefore\ 2\leq d<n\ \land\ \operatorname{mod}(n,d)=0.
\quad \blacksquare\ \text{[Q.E.D.]}
\end{equation*}

\textbf{Smallest divisor is at most sqrt(n).} When $n$ is composite with
smallest divisor $d$, the factor $q = n / d$ satisfies $q \ge d$. Then
$d \cdot d \le d \cdot q = n$, so $d^2 \le n$. This means the scan only
needs to check divisors up to $\sqrt{n}$ --- any divisor beyond that would
have a co-factor below $d$, violating minimality.

\begin{equation*}
\begin{aligned}
n > 1 \;\land\; \neg \operatorname{isPrime}(n) &\Rightarrow \\
d = \operatorname{findSmallestDivisor}(n, 2) &: d \cdot d \leq n.
\end{aligned}
\end{equation*}

\textbf{Proof.} Put $q=n/d$. Since $d$ divides $n$, $dq=n$. Because
$d<n$, $q>1$; hence $q\ge2$. If $q<d$, then $q\in[2,d)$ is a divisor of
$n$, contradicting the minimality of $d$. Hence $d\leq q$, and

\begin{equation*}
d^2\leq dq=n.\quad\blacksquare\ \text{[Q.E.D.]}
\end{equation*}

\textbf{Packaged composite divisor.} The wrapper
\texttt{assertCompositeSmallestPrimeDivisor} combines the previous results
into a reusable form: every composite number has a non-trivial prime
divisor, the divisor really divides the number, and it lies at or below
the square root bound.

\begin{equation*}
\begin{aligned}
n > 1 \;\land\; \neg \operatorname{isPrime}(n)
&\Rightarrow \exists d: \\
&2 \le d < n
\;\land\; \operatorname{isPrime}(d)
\;\land\; d^2 \le n \\
&\qquad \;\land\; \operatorname{Calc.mod}(n,d)=0
  &&\text{[Composite Smallest Prime Divisor]}
\end{aligned}
\end{equation*}

\textbf{Proof.} From the composite assumption,
\texttt{assertCompositeHasDivisorStrictlyBelowN(n)} gives $d < n$ with
$\operatorname{mod}(n, d) = 0$. Let $q = n / d$, so $q \cdot d = n$. If
$q < d$, then $q$ is a divisor of $n$ smaller than $d$, contradicting $d$
being the smallest divisor. Therefore $q \ge d$, and
$d \cdot d \le d \cdot q = n$.

{\raggedright
These properties are verified in the
\href{https://github.com/thiagomata/prime-numbers/blob/master/src/main/scala/v1/chapter5/prime/properties/PrimeProperties.scala}{\texttt{PrimeProperties::assertSmallestDivisorAtMostSqrt}},
\href{https://github.com/thiagomata/prime-numbers/blob/master/src/main/scala/v1/chapter5/prime/properties/PrimeProperties.scala}{\texttt{PrimeProperties::assertCompositeHasDivisorStrictlyBelowN}},
and
\href{https://github.com/thiagomata/prime-numbers/blob/master/src/main/scala/v1/chapter5/prime/properties/PrimeProperties.scala}{\texttt{PrimeProperties::assertCompositeSmallestPrimeDivisor}}.
\par}

\subsection{Finite-Prefix Primality Criterion}
\label{sec:finite-prefix-primality}

The next supporting result is a local primality criterion. If a candidate
number is coprime to all primes in a finite filter list, and every integer
in the range $[2, head)$ has a prime factor among those filters, then the
candidate itself is prime.

This is not Euclid's infinitude theorem; it is a finite-prefix primality
criterion. It turns coverage of all smaller possible divisors into
primality of the candidate.

\begin{equation*}
\begin{aligned}
head > 1
\;\land\; \operatorname{isCoprime}(head,\overline P)
\;\land\;
\forall d\in[2,head),\ \neg \operatorname{isCoprime}(d,\overline P)
&\Rightarrow \operatorname{isPrime}(head).
\end{aligned}
\end{equation*}

The proof is by contradiction over possible divisors. If a divisor $d$ of
$head$ existed in $[2, head)$, the range-coverage assumption would provide
a prime factor from the finite filter list dividing $d$. Divisibility
would then propagate from that factor through $d$ into $head$,
contradicting that $head$ is coprime to every filter prime.

This property is verified in
\href{https://github.com/thiagomata/prime-numbers/blob/master/src/main/scala/v1/chapter5/prime/properties/PrimeProperties.scala}{\texttt{PrimeProperties::assertHeadIsPrime}}.
Its main range helper is
\href{https://github.com/thiagomata/prime-numbers/blob/master/src/main/scala/v1/chapter5/prime/properties/PrimeProperties.scala}{\texttt{PrimeProperties::assertNoDivisorInRangeFromHelper}}.

\subsection{B\'ezout and Prime-Product Lemmas}
\label{sec:bezout-prime-product}

Several arguments about prime-filtered products use a product form of
primality: for a prime $p$, divisibility of a nonnegative product by $p$
can be pushed onto a factor, and if neither nonnegative factor is
divisible by $p$, then the product is not divisible by $p$. The verified
proof goes through B\'ezout's identity.

First, if $0 < h < p$, $p$ is prime, and $h$ is not divisible by $p$, then
$h$ and $p$ have greatest common divisor $1$, and the extended Euclidean
algorithm exposes a linear combination:

\begin{equation*}
\begin{aligned}
\operatorname{isPrime}(p)
\land 0 < h < p
\land \operatorname{mod}(h,p)\ne0
&\Rightarrow
\exists x,y,\ h x + p y = 1.
\end{aligned}
\end{equation*}

For completeness, the subtractive extended Euclidean algorithm repeatedly
replaces the larger positive input by its difference with the smaller.
Common divisors are preserved and the positive sum decreases, so it
reaches equal values. Reversing either subtraction step preserves a
linear-combination witness:

\begin{equation*}
\begin{aligned}
(a-b)x+by=g &\implies ax+b(y-x)=g, \\
ax+(b-a)y=g &\implies a(x-y)+by=g.
\end{aligned}
\end{equation*}

Every common divisor of $h$ and the prime $p$ is either $1$ or $p$; it
cannot be $p$ because $0<h<p$. The terminal gcd is therefore $1$, which
gives the displayed B\'ezout identity. $\blacksquare\ \text{[Q.E.D.]}$

Multiplying that identity by $k$ gives $k h x + k p y = k$. If $p$ divides
$k h$, then $p$ divides both terms on the left and therefore divides $k$.

\begin{equation*}
\begin{aligned}
\operatorname{isPrime}(p)
\land k\ge0
\land h\ge0
\land \operatorname{mod}(h,p)\ne0
\land \operatorname{mod}(kh,p)=0
&\Rightarrow \operatorname{mod}(k,p)=0.
\end{aligned}
\end{equation*}

The contrapositive form used by product and density arguments is:

\begin{equation*}
\begin{aligned}
\operatorname{isPrime}(p)
\land k\ge0
\land h\ge0
\land \operatorname{mod}(k,p)\ne0
\land \operatorname{mod}(h,p)\ne0
&\Rightarrow \operatorname{mod}(kh,p)\ne0.
\end{aligned}
\end{equation*}

{\raggedright
The full proof bodies are verified in
\href{https://github.com/thiagomata/prime-numbers/blob/master/src/main/scala/v1/chapter5/prime/BezoutUtils.scala}{\texttt{BezoutUtils::assertCoprimeLinearCombinationOne}},
\href{https://github.com/thiagomata/prime-numbers/blob/master/src/main/scala/v1/chapter5/prime/BezoutUtils.scala}{\texttt{BezoutUtils::assertPrimeDivKhImpliesDivK}},
and
\href{https://github.com/thiagomata/prime-numbers/blob/master/src/main/scala/v1/chapter5/prime/BezoutUtils.scala}{\texttt{BezoutUtils::assertPrimeProductNotDivisible}}.
\par}
